% Generated draft: requires mathematical and source review.
\documentclass[11pt]{article}
\usepackage[margin=25mm]{geometry}
\usepackage{fontspec}
\setmainfont{Latin Modern Roman}
\usepackage{amsmath,amssymb,mathtools,unicode-math}
\setmathfont{Latin Modern Math}
\usepackage{enumitem,needspace,xurl,hyperref}
\setlength{\parindent}{0pt}
\setlength{\parskip}{5pt}
\setlength{\emergencystretch}{3em}
\newcommand{\N}{\mathbb{N}}
\newcommand{\Z}{\mathbb{Z}}
\newcommand{\Q}{\mathbb{Q}}
\newcommand{\R}{\mathbb{R}}
\newcommand{\C}{\mathbb{C}}
\newcommand{\F}{\mathbb{F}}
\newcommand{\A}{\mathbb{A}}
\newcommand{\PP}{\mathbb{P}}
\newcommand{\E}{\mathbb{E}}
\newcommand{\eps}{\varepsilon}
\newcommand{\dd}{\,\mathrm{d}}
\begin{document}
\section*{UnsolvedMath ID: 30003312}
\textbf{Source problem code:} catalogue entry
\subsection{Definitions and mathematical statement}
The supplied record groups several questions under permanent UnsolvedMath ID \(30003312\), source code \(\text{catalogue entry}\). The formulations below retain the questions separately; some require additional specifications before they determine unique mathematical assertions.

Let \(\mathbb T\) denote the logarithmic-exponential series field named in the record, and let \(\mathbf{No}\) denote the ordered field of surreal numbers, which is a proper class. Write \(\partial_{\mathbb T},\exp_{\mathbb T}\) and \(\partial_{\mathbf{No}},\exp_{\mathbf{No}}\) for the indicated derivations and exponential maps. A derivation on a field \(K\) is an additive map \(\partial:K\to K\) satisfying \(\partial(ab)=a\partial(b)+b\partial(a)\). The record does not specify the construction of \(\mathbb T\) or identify the intended derivation on \(\mathbf{No}\); these choices must be fixed for the following questions.

\begin{enumerate}
\item \textbf{Elementary equivalence.} Under the usual ordered differential-exponential field interpretation, take
\[
\mathcal L=\{0,1,+,-,\times,<,\partial,\exp\}.
\]
Is
\[
(\mathbb T;0,1,+,-,\times,<,\partial_{\mathbb T},\exp_{\mathbb T})
\equiv_{\mathcal L}
(\mathbf{No};0,1,+,-,\times,<,\partial_{\mathbf{No}},\exp_{\mathbf{No}})?
\]
Here elementary equivalence asks whether every first-order \(\mathcal L\)-sentence has the same truth value in the two structures. The record does not explicitly fix the language or any named constants beyond field operations. Since \(\mathbf{No}\) is a proper class, a foundational interpretation of first-order satisfaction in this class structure is also required; the display is a question schema under such an interpretation.

\item \textbf{Exponential-compatible Hardy-field embeddings.} A Hardy field is a field of germs at \(+\infty\) of real-valued differentiable functions, closed under differentiation, with the order given by eventual comparison. Two representatives define the same germ when they agree on some sufficiently far right-hand tail. Such a field \(H\) is closed under exponentiation if \([h]\in H\) implies \([\exp\circ h]\in H\). Its derivation is \(\partial_H[h]=[h']\), and its exponential map is \(\exp_H[h]=[\exp\circ h]\).

Does every Hardy field \(H\) closed under exponentiation admit an injective ordered field homomorphism
\[
\iota:H\longrightarrow\mathbf{No}
\]
for which, for every \(h\in H\),
\[
\iota(\partial_H h)=\partial_{\mathbf{No}}(\iota(h)),
\qquad
\iota(\exp_H h)=\exp_{\mathbf{No}}(\iota(h))?
\]
The explicit original statement identifies \(\mathbf{No}\) as the target. No condition that \(H\) contain all real constant germs, or that \(\iota\) fix all real constants, is supplied. Such conditions must not be added without clarification.

\item \textbf{Trigonometric functions and complex exponentiation.} Can one define sine and cosine on all of \(\mathbf{No}\), and extend exponentiation to
\[
\mathbf{No}[i]=\{a+bi:a,b\in\mathbf{No},\ i^2=-1\}?
\]
For orientation, an extension of real exponentiation as an exponential map would at least be a map \(E:\mathbf{No}[i]\to\mathbf{No}[i]^{\times}\) satisfying
\[
E(z+w)=E(z)E(w),\qquad E(a)=\exp_{\mathbf{No}}(a)
\quad(a\in\mathbf{No}).
\]
These are only minimal extension conditions, not a complete specification of the source question. The record does not state the required real restrictions of sine and cosine, addition formulas, compatibility with differentiation or complex exponentiation, regularity, or meaning of ``define'' (construction versus definability in a specified language). Therefore this part remains a request for suitable structures with unspecified additional properties.

\item \textbf{Expanded structure and an intermediate-value property.} The record asks about good model-theoretic properties, including an intermediate-value property, of
\[
\bigl(\mathbf{No},+,\times,\exp,\widehat f\bigr)_{f\in\mathbb R\langle\omega\rangle}.
\]
Neither \(\mathbb R\langle\omega\rangle\) nor the operation \(f\mapsto\widehat f\), including the arities and domains of the resulting functions, is defined in the supplied record. The class of functions to which the intermediate-value property should apply is likewise unspecified. For a specified unary function \(g\) defined on an interval \(I\subseteq\mathbf{No}\), the usual intermediate-value condition would be
\[
\begin{gathered}
\forall a,b\in I\ \forall y\in\mathbf{No},\\
a<b\ \land\ \min\{g(a),g(b)\}\le y\le\max\{g(a),g(b)\}
\ \Longrightarrow\
\exists c\in I\ \bigl(a\le c\le b\ \land\ g(c)=y\bigr).
\end{gathered}
\]
This explains the condition only after \(g\) and its domain are specified. The record does not determine whether the intended quantification ranges over the named functions, terms, continuous definable functions, or another family, nor which topology or continuity hypotheses are intended.
\end{enumerate}

The background assessment additionally mentions composition and \(p\)-adic questions, but supplies no mathematical formulations for them. They cannot be reconstructed from this record. The supplied label is ``partially solved''; its dated assessment distinguishes a claimed differential Hardy-field embedding result from the exponential-compatible question above. These are supplied curation claims, not independently verified conclusions of this draft.
\subsection{Short English statement}
This record asks whether the indicated series field and surreal field have the same first-order theory with derivation and exponentiation, whether every exponential-closed Hardy field embeds into the surreal field preserving those operations, and whether suitable surreal trigonometric, complex exponential, and intermediate-value structures exist. Several choices and function families are unspecified, so parts remain conditional formulations requiring human mathematical review. The supplied ``partially solved'' assessment is dated curation, not an independently verified status.
\subsection{Sources}
\begin{itemize}
\item Local UnsolvedMath registry: supplied complete record, permanent ID \(30003312\), title ``Model Theory and Analysis of Surreal Exponential Fields,'' source code \(\text{catalogue entry}\). Project registry references: \texttt{data/unsolvedmath/problems.json} and \texttt{catalog/problems.json}. No individual proposer is named in the record.
\item Supplied report citation: \emph{Mini-Workshop: Surreal Numbers, Surreal Analysis, Hahn Fields and Derivations} (2017), \emph{Oberwolfach Reports}, pp.~3313--3372, DOI \texttt{10.4171/owr/2016/60}. \url{https://doi.org/10.4171/owr/2016/60}. The record attributes the questions to workshop open-problem discussions, lists proposed year 2016, and supplies no page pinpoint for this particular entry.
\item Supplied literature source: Aschenbrenner--van den Dries--van der Hoeven. \url{https://doi.org/10.4171/jems/858}. No article title or theorem pinpoint is supplied.
\item Supplied literature source: original OWR report. \url{https://doi.org/10.4171/owr/2016/60}.
\item Supplied literature source: Berarducci--Mantova. \url{https://doi.org/10.4171/jems/769}. No article title or theorem pinpoint is supplied.
\item Literature assessment date, copied exactly: \texttt{2026-08-21}. This is supplied dated curation, not independently verified. Its status is ``partially solved.'' It states that later work proves the differential Hardy-field embedding into \(\mathbf{No}\), while exponential-compatible embedding and additional composition, trigonometric, \(p\)-adic, and model-theoretic questions remain unresolved. The background explicitly calls this dated literature triage and requests independent verification of time-sensitive status and cited claims.
\item The supplied statement-status label is \texttt{corrected\_verified}; its accompanying note says the official report passage was checked and rewritten as a self-contained problem preserving its mathematical content. This is a supplied verification claim. No source was independently consulted for this draft.
\end{itemize}
\end{document}
